\section*{Devoir maison}


\begin{lstlisting}[language=Python]
def recherche(d, x):
    for k, v in d.items() :
        if k == x:
            return v
    return -1
\end{lstlisting}

\begin{Exercise}[title={Recherche linéaire}]

\Question Quel est le type du paramètre d ? Justifiez votre réponse. \\
\emptybox{1cm}

\Question Compléter le tableau d'avancement des variables pour un appel de la fonction recherche(\{9: 2, 7: 1, 11: 11, 6: 7, 0: 2\}, 6) \\

\begin{center}
\large
\begin{tabular}{c|c|c|c|}
\textbf{Tour} & \textbf{k} & \textbf{v} & \textbf{k == x}  \\
\hline
0 &  &  &  \\
1 &  &  &  \\
2 &  &  &  \\
3 &  &  &  \\
4 &  &  &  \\
\end{tabular}
\end{center}

\Question Quelles sont les valeurs de retours de la fonction recherche pour les appels suivants : \\
\begin{itemize}
    \item recherche(\{9: 2, 7: 1, 11: 11\}, 7)
    \item recherche(\{\}, 3)
    \item recherche(\{3: 1, 1: 0\}, 5)
    \item recherche(\{5: 6, 7: 1, 11: -1\}, 11)
\end{itemize}

\Question On remarque que les appels : recherche(\{3: 1, 1: 0\}, 5) et recherche(\{5: 6, 7: 1, 11: -1\}, 11) retournent le même résultat. Expliquez pourquoi cela n'est pas satisfaisant ? \\
\emptybox{5cm}

\Question Proposez une solution pour palier à ce problème.\\
\emptybox{5cm}

\end{Exercise}
